% CAS (Heimdall) -- system functional description: oxygen system.
% Structure after Cooling, Software Engineering for Real-Time Systems (2020), sec. 14.3 [COO].
% Build: latexmk -pdf -interaction=nonstopmode "CAS FD Oxygen.tex"
\documentclass[11pt,a4paper]{article}
% amsmath defines \tag; casfd.sty defines its own \tag for instrument tags. Load amsmath first and
\usepackage{amsmath}
\usepackage{casfd}
\usepackage{adjustbox}

\newcommand{\Otwo}{O\textsubscript{2}}
\newcommand{\degC}{\,\textdegree C}
\newcommand{\TBD}{\textcolor{red}{TBD}}
\setlength{\emergencystretch}{2em}

\casfdheader{Oxygen system}{CAS-FD-20}{A1}

\begin{document}
\maketitle

\begin{abstract}
\noindent This document describes how the oxygen system of the CAS closed pen at Heimdall works,
and how the control system runs it. It is written for the project team, the operators, the
commissioning engineers and the oxygen supplier Normex. It covers every instrument and actuator
of the oxygen system, the control principle, the start and stop sequences, the interlocks, the
alarms and the fault handling, as implemented in the PLC program of 2026-10-21 and as required by
the Norwegian functional description [FD]. Items that are not decided are marked \open{n} and listed
in Section~\ref{sec:open}; decisions are marked \decided{n}.
\end{abstract}

\tableofcontents

% =============================================================================================
\section{Introduction}
\label{sec:intro}

\subsection{Purpose}
Fish in the closed pen get all their oxygen from the water pumped in. The intake water does not
hold enough oxygen for the biomass, so pure oxygen is dissolved into a side stream on each of the
three inlets, and the dissolved \Otwo{} in the pen is controlled to a setpoint. This document is
the \emph{system functional specification} of that oxygen system in the sense of Cooling
[COO, \S14.3]: it describes the system as blocks (Section~\ref{sec:block}), as functions
(Section~\ref{sec:func}), and states the functional (Section~\ref{sec:freq}) and non-functional
(Section~\ref{sec:nfreq}) requirements.

\subsection{Scope}
In scope:
\begin{itemize}[nosep]
  \item the oxygen gas supply and distribution panel up to each injector (supplier Normex), as far
        as it is seen by the control system;
  \item per inlet line $n = 1,2,3$: the side-stream (oxy) pump with its variable speed drive, the
        pressure transmitters, the \Otwo{} injector (oxygenator), the \Otwo{} flow meter and the
        \Otwo{} control valve;
  \item the pen \Otwo{} measurement (Orbit 864 multisensors) and the pen \Otwo{} controller OIC01;
  \item the PLC functions, alarms, interlocks and operator functions for all of the above.
\end{itemize}
Not in scope: the inlet pumps (Framo, see the inlet FD), the emergency oxygenation ([FD] \S3.2.5,
a separate system without connection to the control system), and the CO\textsubscript{2}/TDGP
measurement except where it acts on the oxygen system.

\subsection{Sources}
\begin{small}
\begin{longtable}{@{}P{14mm}P{98mm}P{45mm}@{}}
\caption{Sources. Codes in brackets are used in the text.}\label{tab:sources}\\
\toprule
Code & Document & Used for \\
\midrule
\endfirsthead
\toprule
Code & Document & Used for \\
\midrule
\endhead
\bottomrule
\endfoot
{[COO]} & J.~Cooling, \emph{Software Engineering for Real-Time Systems}, 2020, \S14.3 System functional specifications & Structure of this document \\
{[FD]} & P-15169-CAS-FD-01 Funksjonsbeskrivelse CAS rev~B (Norwegian), \S3.2.4 Oksygentilførsel, tables 12--15, 17 & Required functions \\
{[AL]} & P-15169-CAS-FD-02 Alarmliste CAS Heimdall rev~A & Required alarms \\
{[PID]} & CAS P\&ID and loop diagrams (workfolder-latex), sheets CAS-PID-200, -300, -400, CAS-LAY-700, CAS-LD-802..804, CAS-PID-900/901, -950 & Tags, I/O, open points \\
{[PLC]} & Heimdall PLC project (CODESYS export of 2026-10-21): \texttt{FB.OxySystem.OxyControl}, \texttt{OxyPumpATV650}, \texttt{OxySensors}, \texttt{FB.Inlet.InletSupervisor}, \texttt{FB.Pen.PenControl}, \texttt{FB.RIO.InletRIO}, DUT \texttt{System.OxySystem}, \texttt{PFC200\_Sim} \texttt{ConfigValues} & Implemented behaviour, defaults \\
{[PD]} & CAS program design (workfolder-latex), oxygen dataflow O1--O10 & Software structure \\
{[DB]} & CAS design basis (workfolder-latex) and CasDesignBasis v2.3.12 & Process model, tuning \\
{[O2S]} & Beregningsark dimensjonering ved bortfall av pumper eller strøm, 17.03.26 & \Otwo{} demand, saturation targets \\
{[CT]} & Chucaotech: Oppsett innløsere CT.xlsx; brochure & Injector capacity \\
{[NX]} & Flygt NX 3153 LT 3\textasciitilde{} 411 datasheet, 06.01.2026; MiniCAS II & Side-stream pump \\
{[HZ]} & Risikovurdering -- HAZID -- Oksygen (copy 2026-02-04) & Hazards \\
{[SO]} & Sensor Overview CAS V1.xlsx & Sensors \\
{[ORB]} & Orbit 864 assembly manual DOC100486-04; RDO Core IC spec sheet; RDO Core Modbus interface spec & Pen \Otwo{} sensors \\
{[EAV]} & Schneider ATV600 communication parameters EAV64332 V4.6; ATV6xx programming manual EAV64318 & Drive registers, fault codes \\
{[NMX]} & Normex: revised request sketch 15.10.26; CS Instruments VA~500 datasheet & Oxygen panel \\
\end{longtable}
\end{small}

\subsection{Abbreviations}
\begin{small}
\begin{tabularx}{\textwidth}{@{}P{22mm}L P{22mm}L@{}}
\toprule
AI / AO / DI & analog input / analog output / digital input & OIC & oxygen indicating controller \\
DRIVECOM & CiA 402 drive state machine used by the ATV & PIC & pressure indicating controller \\
ETA, CMD & drive status / command word (Modbus registers) & PI & proportional--integral (controller) \\
FD & functional description & PLC & programmable logic controller (WAGO PFC200) \\
FIC & flow indicating controller & RIO & remote I/O node (WAGO 750-362 coupler) \\
HMI & human--machine interface & SP / PV & setpoint / process value \\
LOX & liquid oxygen & TDGP & total dissolved gas pressure \\
MX & mixer: here the \Otwo{} injector (oxygenator) & VSD & variable speed drive (ATV650) \\
\bottomrule
\end{tabularx}
\end{small}

\subsection{Revision note}
\begin{small}
\begin{tabularx}{\textwidth}{@{}P{10mm}P{22mm}L@{}}
\toprule
Rev. & Date & Change \\
\midrule
A & 2026-10-21 & First issue. Written from [FD] rev~B, the P\&ID set and the PLC program of 2026-10-21. Reflects the choice of Normex as oxygen supplier\decided{1}. \\
A1 & 2026-10-22 & RIO link loss: the line's pressure and flow values freeze with good quality (fault table, OI28). Inlet side: the inlet pumps' speed now follows the shared total of LIC01 and the PLC chooses the number of inlet pumps (CAS-FD-10 rev~B); a line stops when its inlet pump is stopped by the pump count or rotation and starts again when its pump is back in REMOTE. Oxygenation is switched for the whole pen (Pen Oxygen Start / Stop) and each line has a retained Enable, replacing the per-line \var{Command.EnableOxyControl}; valve auto / manual as pulses. \\
\bottomrule
\end{tabularx}
\end{small}

\textbf{Decisions referred to in this document.}
\decided{1} 2026-10-21: Normex is the supplier of the oxygen supply and distribution panel (replaces
the Linde SOLVOX cabinet drawn on CAS-PID-200).
\decided{2} The side-stream pump drive is a Schneider ATV650D15N4 in drive cabinet DC$n$ (drawing
874.004), controlled over Modbus TCP on its embedded Ethernet port. [PD] and [DB] still name an
ATV320 in places; the PLC block is \var{OxyPumpATV650}.

% =============================================================================================
\section{Block diagram description}
\label{sec:block}

Cooling's block diagram description [COO \S14.3.2] lists the physical parts of the system and how
they are connected. Figure~\ref{fig:block} shows the oxygen system as blocks and network links;
Figure~\ref{fig:process} shows one inlet line as a simplified process diagram with the P\&ID tags.

\subsection{Process in one paragraph}
Each of the three Framo inlet pumps (PM-31..33) pushes sea water into the pen through its inlet line.
A small side stream (about 300\,m\textsuperscript{3}/h, 2--4\,\% of the inlet flow [DB]) is
taken from each inlet line by a side-stream pump (oxy pump PU4$n$), pressurised, and led through a
Chucaotech \Otwo{} injector (MX-4$n$) where pure oxygen gas is added and dissolved. The
oxygen-rich side stream returns to the same inlet line and enters the pen with the main flow.
The control system measures the dissolved \Otwo{} in the pen and sets how much oxygen gas each line
gets. A line can only add oxygen while its inlet pump runs: an injector on a stopped inlet cannot
bring oxygen into the pen [DB].

\subsection{Blocks}

\textbf{Oxygen supply and distribution panel (Normex)\decided{1}.} Oxygen gas comes from a central
supply (type \TBD: LOX tank, generator, or the customer's supply line\open{1}) to a distribution panel.
According to the Normex sketch [NMX] the panel has a DN50 header with manometer PI-01.4
($-1$..+6\,bar), and per line a ball valve BV-01.$n$, a flow meter FM-01.$n$, a motorised control
valve MV-01.$n$, a pressure transmitter PT-01.$n$ ($-1$..+6\,bar), a manometer PI-01.$n$ and an
outlet valve OV-02.$n$, with a manual bypass (OV-01.$n$, BV-01.$n$) around each line and cross-over
valves BV-01.7/01.8. The pneumatic shut-off valves PV-01.$n$ and the water-in-pipe transmitter
WD-01.1 are taken out of the scope. The supply pressure is marked ``2.5\,bar ?'' on the sketch\open{1}.
In this document the flow meter and the control valve keep the CAS tags of the P\&ID: FM-01.$n$ =
\itag{FT4$n$}, MV-01.$n$ = \itag{OV4$n$}. The oxygen supply pressure seen by the PLC is \itag{PT40}.
The interface (signals, ranges, valve type) is not yet agreed with Normex\open{2}\textsuperscript{,}\open{3}\textsuperscript{,}\open{8}.

\textbf{Side-stream pump PU4$n$.} Flygt NX 3153 LT, 13.5\,kW, 400\,V, 28\,A, 1455\,rpm, impeller
240\,mm [NX]. Duty point at 50\,Hz: 5.05\,m\textsuperscript{3}/min at 10.2\,m head; shut-off
head 18.3\,m. Leakage and motor temperature sensors are standard; a Flygt MiniCAS II relay in
drive cabinet DC$n$ gives dry contacts \itag{XS4$n$} (leakage) and \itag{TSH4$n$} (temperature).

\textbf{Variable speed drive -FC01 in DC$n$.} Schneider ATV650D15N4, 15\,kW\decided{2}. The PLC
reaches it on the drive's embedded Ethernet port with Modbus TCP. The connection is opened by the
line's \var{InletRIO[i]} function block, which runs a second, separate \var{SCL.sModbusMasterTCP}
master for the drive (the drive has its own IP address). Drive cabinet drawing 874.004 shows no
Ethernet connection yet; an RJ45 DIN-rail socket is proposed\open{20}.

\textbf{Pressure transmitters.} \itag{PT41..43} after the pump (= before the injector) and
\itag{PT44..46} after the injector, 4--20\,mA 2-wire on the line's InletRIO 750-455. A transmitter
\emph{before} the pump is not on the P\&ID; the sensor overview [SO] lists six transmitters
``before and after'' the injector, with make and range not yet chosen\open{17}.

\textbf{\Otwo{} injector MX-4$n$.} Chucaotech CT10 per line. Capacity [CT]: 405\,l/min =
34.7\,kg/h \Otwo{} at 5\,\% gas/liquid ratio (v/v), 1215\,l/min = 104.2\,kg/h at 15\,\%. The PLC
caps the flow of one line at \var{FlowMax} = 104.2\,kg/h, i.e.\ assumes 15\,\%\open{7}.

\textbf{\Otwo{} flow meter FT4$n$ and control valve OV4$n$.} Both in the Normex panel and
both wired to the CommonRIO (Section~\ref{sec:io}). The flow meter in the Normex offer is a CS
Instruments VA~500 thermal mass flow meter (4--20\,mA, not isolated, and Modbus RTU; unit
selectable, e.g.\ kg/h or Nl/min) [NMX]. The PLC expects 4--20\,mA = 0--150\,kg/h\open{3}. The
PLC drives the valve with 4--20\,mA = 0--100\,\% opening; [FD] describes a PWM solenoid valve fed
``via 4--20\,mA'', the Linde cabinet had a Bürkert 2865 proportional valve behind a PWM driver,
and the Normex sketch shows a motorised valve\open{2}.

\textbf{Pen \Otwo{} sensors.} Six Orbit 864 multisensors (In-Situ RDO Core IC: \Otwo{} in mg/L and
\% saturation, temperature, conductivity/salinity, pressure). Three inside the pen
(\itag{OT01..03}, used for control) and three outside (\itag{OT21..23}, reference) [PID CAS-PID-400].
They share one RS-485 line, Modbus RTU 9600\,8N1, slave addresses 1--6, through an 8-port
isolated hub -MB01 to a WAGO 750-652 serial module in the CommonRIO.

\textbf{Control system.} WAGO PFC200 (750-8212) PLC without local I/O. All field I/O is on
WAGO 750-362 Modbus TCP couplers: the \emph{CommonRIO} (-R1-001, PLC/RIO1 cabinet) and one
\emph{InletRIO} per inlet (-R1-100 in the PLC/RIO1, RIO2 and RIO3 cabinets). The Framo pump
controllers (PCC$n$) are OPC UA servers; the inlet pump state and heartbeat from them give the
permit for the oxygen line. The HMI reads and writes \var{GVLMain} through the PLC's OPC UA server.

\begin{figure}[htbp]
\centering
\begin{adjustbox}{max width=\linewidth}\begin{tikzpicture}[x=1cm, y=1cm, font=\small]
  % network row
  \node[ext, minimum width=22mm] (hmi) at (0.6,0) {HMI\\\footnotesize OPC UA client};
  \node[ext, minimum width=70mm] (sw) at (7.6,0) {Ethernet switch, network cabinet NC};
  \node[ext, minimum width=26mm] (pcc) at (14.4,0) {Framo PCC$n$\\\footnotesize inlet pump PM-3$n$};
  % control row
  \node[proc, minimum width=22mm] (plc) at (0.6,-2.4) {PLC PFC200\\\footnotesize 750-8212};
  \node[proc, minimum width=30mm] (crio) at (4.6,-2.4) {CommonRIO\\\footnotesize 750-362 -R1-001};
  \node[proc, minimum width=30mm] (irio) at (9.0,-2.4) {InletRIO $n$\\\footnotesize 750-362 -R1-100};
  \node[proc, minimum width=26mm] (vsd) at (14.4,-2.4) {ATV650 -FC01\\\footnotesize drive cabinet DC$n$};
  % field row
  \node[ext, minimum width=26mm, text width=25mm] (orb) at (0.9,-5.2) {750-652 + hub\\Orbit 864 $\times$6\\\footnotesize OT01..03, OT21..23};
  \node[ext, minimum width=34mm, text width=33mm] (nmx) at (5.0,-5.2) {Normex \Otwo{} panel\\\footnotesize FT41..43, OV41..43, PT40};
  \node[ext, minimum width=30mm, text width=29mm] (pts) at (9.0,-5.2) {Line $n$ field\\\footnotesize PT4$n$, PT4($n$+3), XS/TSH4$n$};
  \node[ext, minimum width=26mm] (pu) at (14.4,-5.2) {PU4$n$\\\footnotesize NX 3153, 13.5 kW};
  % links
  \draw[very thick] (hmi) -- (sw) node[lab, pos=0.5, above=1pt] {OPC UA};
  \draw[very thick] (sw) -- (pcc) node[lab, pos=0.5, above=1pt] {OPC UA};
  \draw[very thick] (plc.north) -- (plc.north |- sw.south);
  \draw[very thick] (crio.north) -- (crio.north |- sw.south);
  \draw[very thick] (irio.north) -- (irio.north |- sw.south);
  \draw[very thick] (vsd.north) -- (vsd.north |- sw.south);
  \node[lab] at (11.9,-1.0) {Modbus TCP};
  \draw[sig] (crio.south) ++(-0.6,0) -- ++(0,-1.0) -| (orb.north);
  \node[lab] at (2.2,-3.95) {RS-485, Modbus RTU};
  \draw[sig, <->] (crio.south) ++(0.4,0) -- ++(0,-1.0) -| (nmx.north);
  \node[lab] at (6.0,-3.95) {4--20 mA AI/AO};
  \draw[sig] (pts.north) -- (irio.south) node[lab, pos=0.5] {4--20 mA, DI};
  \draw[very thick] (vsd.south) -- (pu.north) node[lab, pos=0.5] {motor cable};
  \draw[sig] (pts.east) -- (pu.west) node[lab, pos=0.5, above=1pt] {XS/TSH};
  \draw[sig] (vsd.west) -- (irio.east) node[lab, pos=0.5, above=1pt] {2nd master};
\end{tikzpicture}\end{adjustbox}
\caption{Block diagram of the oxygen system for one line $n$ (three lines in total; CommonRIO,
PLC, HMI and the Orbit sensors are shared). Thick: Ethernet; dashed: field signals. The drive
is a Modbus TCP slave of the PLC; the ``2nd master'' runs in the line's \texttt{InletRIO} block.}
\label{fig:block}
\end{figure}

\begin{figure}[htbp]
\centering
\begin{adjustbox}{max width=\linewidth}\begin{tikzpicture}[x=1cm, y=1cm]
  % main water line and pen
  \draw[pipe, draw=blue!70!black] (0,0) -- (13.9,0);
  \node[font=\scriptsize, anchor=north west, text=blue!70!black] at (0,-0.1) {from inlet pump PM-31 (LW-0041)};
  \node[draw, rounded corners=8pt, minimum width=22mm, minimum height=30mm, fill=blue!4,
        font=\small\bfseries] (pen) at (15.1,0.2) {};
  \node[font=\small\bfseries] at (15.1,-1.0) {PEN01};
  % side stream
  \draw[very thick, draw=blue!70!black] (1.4,0) -- (1.4,1.5) -- (2.15,1.5);
  \node[draw, circle, minimum size=7mm, fill=white] (pu) at (2.5,1.5) {};
  \draw (pu.north) -- (pu.east) -- (pu.south);
  \node[font=\scriptsize, anchor=north] at (2.5,1.05) {PU41};
  \draw[very thick, draw=blue!70!black] (pu.east) -- (6.1,1.5);
  \node[draw, fill=blue!8, minimum width=14mm, minimum height=8mm, font=\scriptsize, align=center]
       (mx) at (6.8,1.5) {MX-41\\injector};
  \draw[pipe, draw=blue!70!black] (7.5,1.5) -- (10,1.5) -- (10,0.08);
  \node[font=\scriptsize, text=blue!70!black, anchor=south] at (4.6,1.55) {side stream LW-0042};
  % drive
  \node[ext, minimum height=7mm, font=\scriptsize] (vsd) at (2.5,3.0) {ATV650 SC41};
  \draw[densely dotted, thick] (vsd.south) -- (pu.north);
  % PT41, PIC41
  \draw (4.2,1.5) -- (4.2,2.15);
  \node[inst] (pt41) at (4.2,2.6) {PT\\41};
  \node[ctrl] (pic) at (4.2,4.0) {PIC\\41};
  \draw[sig] (pt41) -- (pic);
  \draw[sig] (pic.west) -- (vsd.north east);
  % PT44
  \draw (8.7,1.5) -- (8.7,2.15);
  \node[inst] (pt44) at (8.7,2.6) {PT\\44};
  % O2 supply
  \node[ext, minimum width=24mm, text width=23mm, font=\scriptsize] (sup) at (14.8,4.6) {\Otwo{} supply,\\Normex panel};
  \draw[pipe, draw=teal] (sup.west) -- (6.8,4.6) -- (mx.north);
  \draw (13.1,4.6) -- (13.1,5.15);
  \node[inst] (pt40) at (13.1,5.6) {PT\\40};
  \draw (10.6,4.6) -- (10.6,5.15);
  \node[inst] (ft41) at (10.6,5.6) {FT\\41};
  % control valve (bow tie + actuator)
  \draw[fill=white] (8.3,4.35) -- (8.3,4.85) -- (8.9,4.35) -- (8.9,4.85) -- cycle;
  \draw (8.6,4.6) -- (8.6,5.1);
  \draw[fill=white] (8.35,5.1) rectangle (8.85,5.35);
  \node[font=\scriptsize, anchor=north] at (8.6,4.3) {OV41};
  \node[ctrl] (fic) at (9.6,6.5) {FIC\\41};
  \draw[sig] (ft41) -- (fic);
  \draw[sig] (fic.west) -| (8.6,5.35);
  % pen O2, OIC01, FY01
  \node[inst] (ot) at (15.1,0.75) {OT\\01..03};
  \node[ctrl] (oic) at (15.1,2.6) {OIC\\01};
  \draw[sig] (ot) -- (oic);
  \node[inst, densely dashed] (fy) at (12.3,2.6) {FY\\01};
  \draw[sig] (oic) -- (fy);
  \draw[sig] (fy.north) -- (12.3,6.5) -- (fic.east);
  \node[font=\scriptsize, anchor=west, align=left] at (12.4,6.75) {SP: share of total};
  \node[font=\scriptsize, anchor=west, align=left, text=teal] at (6.9,3.9) {\Otwo{} gas};
\end{tikzpicture}\end{adjustbox}
\caption{Simplified process diagram of oxygen line 1 with the tags of [PID CAS-PID-200/300/400]
(lines 2 and 3 alike: PU42/43, PT42/43, MX-42/43, PT45/46, FT42/43, OV42/43). Blue: water;
green: \Otwo{} gas; dashed: signals. FY01 shares OIC01's total \Otwo{} flow between the lines.}
\label{fig:process}
\end{figure}

\subsection{Typical tags per line}
\begin{center}
\begin{small}
\begin{tabularx}{\textwidth}{@{}P{34mm}P{22mm}P{22mm}P{22mm}L@{}}
\toprule
Item & Line 1 & Line 2 & Line 3 & PLC object \\
\midrule
Side-stream pump, drive & PU41, DC1 -FC01 & PU42, DC2 -FC01 & PU43, DC3 -FC01 & \var{OxySystem[i].OxyPump} \\
Pressure after pump / after injector & PT41 / PT44 & PT42 / PT45 & PT43 / PT46 & \var{.Input.OxyPumpPressure} / \var{.OxynatorOutletPressure} \\
Injector & MX-41 & MX-42 & MX-43 & -- \\
\Otwo{} flow, valve & FT41, OV41 & FT42, OV42 & FT43, OV43 & \var{.Input.OxygenFlow}, \var{.OxyValve} \\
Controllers & PIC41, FIC41 & PIC42, FIC42 & PIC43, FIC43 & \var{OxyControl[i]} \\
RIO node & InletRIO1 & InletRIO2 & InletRIO3 & \var{InletRIO[i]}, $i = n-1$ \\
\bottomrule
\end{tabularx}
\end{small}
\end{center}
[FD] numbers the pressure transmitters PT41/43/45 (in) and PT42/44/46 (out); the cabinet drawings
and this document use PT41--43 after the pump and PT44--46 after the injector\open{27}.

% =============================================================================================
\section{Functional diagram description}
\label{sec:func}

Cooling's functional diagram description [COO \S14.3.3] shows what the system does, as functions
and signal flows, independent of the hardware. The oxygen system has three layers:
\begin{enumerate}[nosep]
  \item \textbf{Pen \Otwo{} control} (one per pen): OIC01 sets the \emph{total} oxygen gas flow
        to the pen, and FY01 shares it between the lines in auto.
  \item \textbf{Line flow control} (one per line): FIC4$n$ sets the valve OV4$n$ so the line's
        measured \Otwo{} flow follows its share.
  \item \textbf{Side-stream pressure control} (one per line): PIC4$n$ sets the pump speed so the
        water pressure after the pump (into the injector) is constant.
\end{enumerate}
Around these sit the line start sequence, the valve permit and the drive state machine.

\begin{figure}[htbp]
\centering
\begin{adjustbox}{max width=\linewidth}\begin{tikzpicture}[x=1cm, y=1cm, font=\small]
  % pen layer
  \node[ext, minimum width=16mm] (sp) at (0,0) {SP\\\footnotesize mg/L};
  \node[sm] (oic) at (2.6,0) {OIC01\\\footnotesize PI};
  \node[sm, text width=26mm] (fy) at (6.0,0) {FY01 share\\\footnotesize $\min\!\big(\frac{T-M}{n_A}, F_\text{max}\big)$};
  \node[sm] (fic) at (9.5,0) {FIC4$n$\\\footnotesize PI};
  \node[sm] (ramp) at (12.0,0) {FY4$n$\\\footnotesize 100\,\%/2\,s};
  \node[sm, minimum width=10mm] (perm) at (14.0,0) {I4\\\footnotesize\&};
  \node[ext] (ov) at (15.8,0) {OV4$n$};
  \draw[flow] (sp) -- (oic);
  \draw[flow] (oic) -- node[lab, above=2pt] {$T$ kg/h} (fy);
  \draw[flow] (fy) -- node[lab, above=2pt] {SP kg/h} (fic);
  \draw[flow] (fic) -- node[lab, above=2pt] {\%} (ramp);
  \draw[flow] (ramp) -- (perm);
  \draw[flow] (perm) -- node[lab, above=2pt] {mA} (ov);
  % feedbacks
  \node[ext] (oy) at (2.6,-1.9) {OY01\\\footnotesize mean OT01..03};
  \draw[flow] (oy) -- node[lab, right] {PV mg/L} (oic);
  \node[ext] (ft) at (9.5,-1.9) {FT4$n$};
  \draw[flow] (ft) -- node[lab, right] {PV kg/h} (fic);
  \node[ext, text width=26mm] (man) at (6.0,-1.9) {lines not in auto\\\footnotesize measured flow $M$};
  \draw[flow] (man) -- (fy);
  \node[ext, text width=22mm] (pv) at (14.0,-1.9) {permit:\\\footnotesize Fig.~\ref{fig:permit}};
  \draw[flow] (pv) -- (perm);
  % pressure layer
  \node[ext, minimum width=16mm] (psp) at (0,-4.0) {SP\\\footnotesize 4.0 bar};
  \node[sm] (pic) at (2.6,-4.0) {PIC4$n$\\\footnotesize PI};
  \node[sm, text width=24mm] (sel) at (6.0,-4.0) {select\\\footnotesize PIC or preset 40 Hz};
  \node[sm] (drv) at (9.5,-4.0) {OxyPump-\\ATV650};
  \node[ext] (atv) at (12.0,-4.0) {ATV650};
  \node[ext] (pu) at (14.0,-4.0) {PU4$n$};
  \node[ext] (pt) at (2.6,-5.8) {PT4$n$ (after pump)};
  \draw[flow] (psp) -- (pic);
  \draw[flow] (pic) -- node[lab, above=2pt] {Hz} (sel);
  \draw[flow] (sel) -- node[lab, above=2pt] {20--50 Hz} (drv);
  \draw[flow] (drv) -- node[lab, below=2pt] {CMD, LFR} (atv);
  \draw[flow] (atv) -- (pu);
  \draw[flow] (pt) -- node[lab, right] {PV bar} (pic);
  \draw[sig] (pu.south) -- ++(0,-0.9) -| node[lab, pos=0.25, below] {water pressure} (pt.east);
\end{tikzpicture}\end{adjustbox}
\caption{Control scheme [COO Fig.~7.8 style]. Top: pen \Otwo{} cascade, OIC01 over FIC4$n$ for each
line in auto. Bottom: pressure loop of each line. $T$ = OIC01 output, $M$ = measured flow of the
active lines not in auto, $n_A$ = number of lines in auto, $F_\text{max}$ = \texttt{FlowMax}.}
\label{fig:cascade}
\end{figure}

\subsection{Pen \texorpdfstring{\Otwo}{O2} measurement (OY01)}
The pen \Otwo{} used for control is the mean dissolved \Otwo{} concentration (mg/L) of the Orbit
sensors that are \emph{selected for control} (\var{Orbit864SensorConfig[k].EnableForAutoControl},
retained; default sensors 1--3 inside the pen), \emph{enabled}, and of good quality. Sensors with
communication failure or bad quality drop out of the mean at once. When none is left, the pen
\Otwo{} measurement is bad (\var{Pen.Status.OxyControlOK} = FALSE). The mean saturation (\%) of the
same sensors is calculated for display. [FD] asks instead for one sensor (OT03) selected by the
operator, and on failure a switch to the lowest of the others; the PLC averages the selected
sensors (deviation, Table~\ref{tab:dev}).

\subsection{Pen oxygen controller OIC01}
OIC01 is a PI controller from the pen \Otwo{} [mg/L] to the \emph{total} oxygen gas flow to the
pen [kg/h]. Because the pen responds to the total flow, its gain does not change with the number
of lines running; one parameter set is enough (the three sets of [FD] table~13 are not needed).
\begin{itemize}[nosep]
  \item \textbf{Enabled} when the pen is started (\var{PenStart}) and at least one line is in auto.
  \item \textbf{Output limits}: 0 to the sum of \var{FlowMax} of the lines in auto plus the
        measured flow of the active lines not in auto (and \var{MaxOutput} if set), so the
        integral cannot wind up beyond what the lines can deliver.
  \item \textbf{Bumpless start}: on the first enabled cycle the output starts at the measured
        total flow.
  \item \textbf{Manual}: the operator can set the total flow by hand
        (\var{OxyController.Command.Manual}, \var{.ManualOutput}).
  \item \textbf{Bad \Otwo{} measurement}: the output is frozen at its last value (interlock I7
        on CAS-PID-400) for as long as the measurement is bad\open{14}.
\end{itemize}

\subsection{Sharing between the lines (FY01)}
A line is \emph{in auto} when oxygenation is active, the valve permit is given, the FIC is
enabled, the valve is not in manual and the flow measurement is good. Every cycle:
\begin{equation*}
  \text{SP}_i = \min\!\Big(F_{\text{max},i},\ \frac{T - \sum_{j\notin A}F_j}{n_A}\Big)
  \quad \text{for each line } i \text{ in auto } (A),
\end{equation*}
where $F_j$ is the measured flow of an active line that is not in auto. A line that drops out of
auto, or a manual line whose valve the operator moves, is made up by the other lines in the
same cycle, not slowly through the integral of OIC01. A line not in auto shows its measured flow
as its setpoint.

\subsection{Line flow controller FIC4\texorpdfstring{$n$}{n} and valve OV4\texorpdfstring{$n$}{n}}
FIC4$n$ is a PI controller from the line's \Otwo{} flow [kg/h] to the valve opening 0--100\,\%.
Its output passes a rate limiter (FY4$n$): the valve command changes at most 100\,\% per
\var{ValveRampTime} (2\,s), computed with the task cycle time. The rate-limited command is written
to the 750-563 as 0--65535 = 4--20\,mA.
\begin{itemize}[nosep]
  \item \textbf{Auto}: SP = the line's share from FY01.
  \item \textbf{Valve manual}: the operator sets the opening (\var{OxyValve.Command.Manual},
        \var{.ManualSetpoint}); the FIC tracks it, so the return to auto is bumpless. The line then
        counts as ``not in auto'' with its measured flow.
  \item \textbf{Flow measurement bad}: the FIC holds the last valve position and the line leaves
        auto; its flow is then counted as 0 in the sharing\open{15}.
  \item \textbf{No permit or oxygenation not active}: output forced to 0\,\% = 4\,mA = valve closed,
        also in valve manual (Section~\ref{sec:permit}).
\end{itemize}

\subsection{Side-stream pressure controller PIC4\texorpdfstring{$n$}{n}}
PIC4$n$ is a PI controller from the pressure after the pump (PT4$n$) to the pump speed, limited
to 20--50\,Hz. Setpoint \var{PressureSetpoint} = 4.0\,bar (absolute is assumed\open{17}). It is
active only while the pump runs. While the pump is starting, in pump manual, with
\var{PressureController.Enabled} = FALSE, or with a bad pressure measurement, the pump runs at
the preset speed (\var{PresetSpeed_Hz} = 40\,Hz) or, in manual, at the operator's speed. The
purpose is a constant water pressure and flow through the injector, which keeps its gas transfer
stable. Whether the pump can reach 4.0\,bar is not confirmed: the datasheet head at the duty
point is 10.2\,m (about 1\,bar) at 50\,Hz\open{18}.

\subsection{Line start and stop sequence}
\label{sec:seq}
Oxygenation is started and stopped for the whole pen (Oxygen Start / Stop, CAS-FD-10 rev~B; replaces
[FD]'s group start per line); each line takes part when it is enabled. Figure~\ref{fig:seq} shows the
sequence.
\begin{enumerate}[nosep]
  \item The operator gives Oxygen Start (\var{Pen.Command.OxygenStart}, accepted while the pen is Starting or
        Running; \var{Pen.Status.OxygenOn}). Pen Stop or Oxygen Stop turns it off.
  \item The line is \emph{started} when oxygenation is on AND the line is enabled (\var{OxyConfig[i].Enable},
        retained) AND the inlet pump of the same line is
        in REMOTE with a live heartbeat (\var{Inlet[i].InletPump.Status.OxyPermit}, interlock
        I5) AND no emergency stop\open{22}\textsuperscript{,}\open{23}.
  \item In pump auto the oxy pump is started; it runs at the preset speed until it reaches
        the minimum speed, then PIC4$n$ takes over.
  \item When the pump has been running for \var{RunDelay} (30\,s), \emph{oxygenation is active}.
  \item With the valve permit (Section~\ref{sec:permit}) the valve may open and the line joins the
        sharing (in auto).
  \item \textbf{Stop}: when the operator stops the line, the inlet permit is lost, or the pump
        stops, oxygenation ends at once and the valve command goes to 4\,mA (closed) in the same
        cycle; the 2\,s ramp applies only while the valve is permitted. The pump stops when the line
        is not started (pump in auto).
\end{enumerate}
In pump manual the operator can start and stop the oxy pump alone (\var{OxyPump.Command.Start/Stop}),
also without the inlet permit; oxygenation still needs the line to be started.

\begin{figure}[htbp]
\centering
\begin{adjustbox}{max width=\linewidth}\begin{tikzpicture}[x=1cm, y=1cm, font=\small]
  \tikzset{st/.style={state, text width=34mm}}
  \node[st] (s0) at (0,0) {Line stopped\\\footnotesize pump off, valve 4 mA};
  \node[st] (s1) at (6.2,0) {Pump starting\\\footnotesize preset speed 40 Hz};
  \node[st] (s2) at (12.4,0) {Pump running\\\footnotesize PIC4$n$ active, valve closed};
  \node[st] (s3) at (12.4,-3.2) {Oxygenation active\\\footnotesize FIC4$n$ in auto};
  \node[st] (s4) at (6.2,-3.2) {Valve held closed\\\footnotesize alarm ValvePressureLow};
  \draw[flow] (s0) -- node[lab, above=3pt] {OxygenOn AND Enable\\AND inlet permit} (s1);
  \draw[flow] (s1) -- node[lab, above=3pt] {speed $\geq$\\MinSpeed $-$ 0.5 Hz} (s2);
  \draw[flow] (s2) -- node[lab, right=3pt] {running for\\RunDelay 30 s} (s3);
  \draw[flow] ([yshift=2mm]s3.west) -- node[lab, above=3pt] {no valve permit} ([yshift=2mm]s4.east);
  \draw[flow] ([yshift=-2mm]s4.east) -- node[lab, below=3pt] {permit again} ([yshift=-2mm]s3.west);
  \draw[flow] (s3.south) -- ++(0,-0.8) -| (s0.south);
  \draw (s4.south) -- ++(0,-0.8);
  \node[lab, anchor=north] at (6.2,-4.75) {stop command, inlet permit lost, pump fault or pump stopped $\rightarrow$ valve closed, pump stopped};
\end{tikzpicture}\end{adjustbox}
\caption{Line start/stop sequence (per line, pump in auto). Any state returns to ``Line stopped''
on the stop conditions.}
\label{fig:seq}
\end{figure}

\subsection{Valve permit (interlock I4)}
\label{sec:permit}
The valve may open only when oxygen can be dissolved: the side stream must flow with pressure.
\begin{figure}[htbp]
\centering
\begin{adjustbox}{max width=\linewidth}\begin{tikzpicture}[x=1cm, y=1cm, font=\small]
  \node[anchor=east] (a1) at (0,1.2) {pump state = Running};
  \node[anchor=east] (a2) at (0,0.4) {PT4$n$ enabled and good};
  \node[anchor=east] (a3) at (0,-0.4) {PT4$n$ $>$ ValveOpenPressure (3.0 bar)};
  \node[draw, minimum width=10mm, minimum height=26mm, fill=yellow!12] (g1) at (1.2,0.4) {\&};
  \foreach \y in {1.2,0.4,-0.4} \draw (0.05,\y) -- (0.7,\y);
  \node[anchor=east] (a4) at (0,-0.9) {oxygenation active};
  \node[draw, minimum width=10mm, minimum height=16mm, fill=yellow!12] (g2) at (4.4,0.1) {\&};
  \draw[flow] (g1.east) -- node[lab, above=1pt] {ValvePermit} (g2.west |- 0,0.4);
  \draw[flow] (0.05,-0.9) -- (3.2,-0.9) -- (3.2,-0.2) -- (g2.west |- 0,-0.2);
  \node[anchor=west, text width=60mm] at (5.3,0.1) {\textbf{ValveOpen}: FIC4$n$ may drive OV4$n$.\\
    Otherwise output 0 = 4 mA (closed), also in valve manual.};
  % alarm
  \node[anchor=east] at (0,-2.2) {oxygenation active};
  \node[anchor=east] at (0,-3.0) {NOT ValvePermit};
  \node[draw, minimum width=10mm, minimum height=14mm, fill=yellow!12] (g3) at (4.4,-2.6) {\&};
  \draw[flow] (0.05,-2.2) -- (g3.west |- 0,-2.2);
  \draw[flow] (0.05,-3.0) -- (g3.west |- 0,-3.0);
  \node[anchor=west, text width=60mm] at (5.3,-2.6) {\textbf{PAL4$n$}: alarm ValvePressureLow\\(operator acknowledge)};
\end{tikzpicture}\end{adjustbox}
\caption{Valve permit I4 and alarm PAL4$n$ (\texttt{OxyControl}).}
\label{fig:permit}
\end{figure}
\var{ValveOpenPressure} (3.0\,bar) is on the same scale as \var{PressureSetpoint} (4.0\,bar); if
the transmitters turn out to read gauge pressure, both must be lowered together by 1\,bar\open{17}.

\subsection{Drive state machine (DRIVECOM)}
The PLC walks the ATV650 through the DRIVECOM (CiA 402) states by writing the command word CMD
(8501) and reading the status word ETA (3201) every 100\,ms (block \var{OxyPumpATV650}).
Speed reference LFR (8502) is written in 0.1\,Hz, limited to 20--50\,Hz. A fault is reset by an
operator reset, which sends \texttt{0x0080} for 0.5\,s. The drive must have its command and
reference channel set to the embedded Ethernet port (Fr1, Cd1).
\begin{figure}[htbp]
\centering
\begin{adjustbox}{max width=\linewidth}\begin{tikzpicture}[x=1cm, y=1cm, font=\small]
  \node[state] (sod) at (0,0) {Switch on disabled\\\footnotesize ETA \texttt{0x40}};
  \node[state] (rts) at (4.8,0) {Ready to switch on\\\footnotesize ETA \texttt{0x21}};
  \node[state] (son) at (9.6,0) {Switched on\\\footnotesize ETA \texttt{0x23}};
  \node[state] (oe) at (9.6,-2.6) {Operation enabled\\\footnotesize ETA \texttt{0x27}: running};
  \node[state] (qs) at (4.8,-2.6) {Quick stop active\\\footnotesize ETA \texttt{0x07}};
  \node[state] (flt) at (0,-2.6) {Fault\\\footnotesize ETA \texttt{0x08} (mask \texttt{0x4F})};
  \draw[flow] (sod) -- node[lab, above=2pt] {CMD \texttt{0x0006}} (rts);
  \draw[flow] (rts) -- node[lab, above=2pt] {Run: \texttt{0x0007}} (son);
  \draw[flow] ([xshift=-3mm]son.south) -- node[lab, left=2pt] {Run: \texttt{0x000F}} ([xshift=-3mm]oe.north);
  \draw[flow] ([xshift=3mm]oe.north) -- node[lab, right=2pt] {stop: \texttt{0x0007} (ramp)} ([xshift=3mm]son.south);
  \draw[flow] (qs) -- node[lab, left=2pt, text width=18mm] {\texttt{0x0006}} (rts);
  \draw[flow] (flt) -- node[lab, left=2pt, text width=24mm] {reset \texttt{0x0080}} (sod);
\end{tikzpicture}\end{adjustbox}
\caption{DRIVECOM states as handled by \texttt{OxyPumpATV650} (status masked with \texttt{0x6F}). Any
other state is answered with \texttt{0x0006}. Quick stop active raises the alarm EmergencyStop.}
\label{fig:drivecom}
\end{figure}

Pump state for the logic: \emph{Unavailable} (no communication or fault), \emph{Available}
(stopped, ready), \emph{Starting} (Run requested, below minimum speed), \emph{Running} (operation
enabled at or above 19.5\,Hz), \emph{Stopping} (ramping down).

% =============================================================================================
\section{Functional requirements}
\label{sec:freq}

This section states the functional requirements in the order of [COO \S14.3.4a]. ``Shall''
statements are requirements; ``PLC:'' says how the program of 2026-10-21 fulfils them.

\subsection{Processes}
\begin{enumerate}[label=F\arabic*, nosep]
  \item The system shall hold the pen \Otwo{} at the operator setpoint by adding oxygen gas to the
        side stream of each running inlet (OIC01, FY01, FIC4$n$).
  \item Each line shall run its side-stream pump at constant pressure into the injector (PIC4$n$).
  \item A line shall add oxygen only while its inlet pump runs under remote control, its side-stream
        pump runs and the pressure after the pump is above the valve-open limit.
  \item The total oxygen flow shall be shared equally between the lines in auto, up to each line's
        maximum; lines in manual shall be accounted for with their measured flow.
  \item The system shall measure and display pen \Otwo{} (mg/L, \% saturation), temperature and
        salinity at six positions, three inside and three outside the pen.
\end{enumerate}

\subsection{Inputs and outputs}
\label{sec:io}
\begin{small}
\begin{longtable}{@{}P{20mm}P{46mm}P{42mm}P{45mm}@{}}
\caption{Inputs and outputs of the oxygen system [PID CAS-PID-900/901].}\label{tab:io}\\
\toprule
Tag & Signal & Node / channel & PLC variable \\
\midrule
\endfirsthead
\toprule
Tag & Signal & Node / channel & PLC variable \\
\midrule
\endhead
\bottomrule
\endfoot
\multicolumn{4}{@{}l}{\textit{Inputs}} \\
PT4$n$ & pressure after pump, 4--20 mA & InletRIO$n$ -R1-106 ch3, word 2 & \var{OxySystem[i].Input.OxyPumpPressure} \\
PT4($n$+3) & pressure after injector, 4--20 mA & InletRIO$n$ -R1-106 ch4, word 3 & \var{OxySystem[i].Input.OxynatorOutletPressure} \\
XS41 / TSH41 & pump leakage / motor temp., DI & InletRIO1 -R1-104 ch8/9, word 5.7/5.8 (line 1 only) & \var{.Input.OxyPumpLeakDetedtion} / \var{.OxyPumpHighTemperature} \\
FT41..43 & \Otwo{} flow, 4--20 mA & CommonRIO -R1-008 ch1--3, words 24--26 & \var{OxySystem[i].Input.OxygenFlow} \\
PT40 & \Otwo{} supply pressure, 4--20 mA & CommonRIO -R1-008 ch4, word 27 & \var{Pen.Input.OxygenInjectionPressure} \\
OT01..03, OT21..23 & Orbit 864, Modbus RTU & CommonRIO -R1-003 750-652, words 0--11 & \var{Pen.Input.Orbit864Sensor[0..5]} \\
SC4$n$ & drive status ETA, RFR, LCR, THD, OPR, LFT & ATV650 DC$n$, Modbus TCP & \var{OxySystem[i].OxyPump.Input} \\
-- & inlet permit (Framo REMOTE + heartbeat) & PCC$n$, OPC UA & \var{Inlet[i].InletPump.Status.OxyPermit} \\
\multicolumn{4}{@{}l}{\textit{Outputs}} \\
OV41..43 & valve opening, 4--20 mA & CommonRIO -R1-014 ch1/2, -R1-015 ch1, out words 20--22 & \var{OxySystem[i].OxyValve.Output.RawPosition} \\
SC4$n$ & CMD, LFR & ATV650 DC$n$, Modbus TCP & \var{OxySystem[i].OxyPump.Output} \\
\end{longtable}
\end{small}

\subsection{Start-up and shut-down}
\begin{enumerate}[label=S\arabic*, nosep]
  \item \textbf{PLC start} (power-up, download): all retained configuration is kept
        (Section~\ref{sec:params}); oxygenation starts off; valves at 4\,mA. The operator gives Oxygen
        Start again. Orbit sensor selection for control is retained (default sensors 1--3).
  \item \textbf{Pen start}: OIC01 is enabled with the pen start (\var{Pen.Command.Start}) once a line
        is in auto, and starts bumplessly from the measured total flow.
  \item \textbf{Line start}: Section~\ref{sec:seq}; [FD] requires auto restart after power return
        (table 2.2: auto-restart ``Yes'', sequence 3, after the inlet pumps). PLC: \var{Pen.Status.OxygenOn}
        is not retained over a power cut; whether oxygenation shall restart by itself is \TBD\open{23}.
  \item \textbf{Line stop}: operator stop, inlet permit lost or pump stopped closes the valve at once
        and stops the pump (in auto).
  \item \textbf{Normex supply start/stop}: local at the panel (manual valves). No PLC-controlled master
        shut-off valve exists in the present Normex scope\open{5}.
\end{enumerate}

\subsection{Fault identification and recovery}
\begin{small}
\begin{longtable}{@{}P{34mm}P{40mm}P{48mm}P{30mm}@{}}
\caption{Faults, their detection and the system's response.}\label{tab:faults}\\
\toprule
Fault & Detection & Response (PLC 2026-10-21) & Recovery \\
\midrule
\endfirsthead
\toprule
Fault & Detection & Response (PLC 2026-10-21) & Recovery \\
\midrule
\endhead
\bottomrule
\endfoot
Drive fault & ETA state ``fault''; code in LFT & Pump Unavailable, line drops out of auto, valve closed, alarms DriveError + PumpCommonError & Operator reset (\texttt{0x0080}), then restart \\
Drive communication lost & \var{DriveMaster.comOk} FALSE (socket timeout 5\,s) & Pump Unavailable, valve closed, DriveCommunicationError. The drive's own reaction is set by ETHL (factory: freewheel stop) & Automatic when the link returns; [FD] wants the pump to keep running at preset speed\open{19} \\
Drive not in remote & ETA bit 9 = 0 & Alarm DriveNotInRemote; commands ignored by the drive & Set drive to remote \\
Pump does not start / stop & not turning 10\,s after Run; turning 10\,s after stop & StartUpTimeout / StopTimeout & Operator ack \\
Pressure after pump low & PT4$n$ $\leq$ 3.0\,bar or bad, while oxygenation active & Valve forced closed (4\,mA), alarm ValvePressureLow; PIC on preset speed if PT bad & Automatic when pressure returns; operator ack \\
Line flow measurement bad & 750-455 diagnostic (wire break, under/over range) & Line leaves auto; valve holds last position; its flow counts as 0 in the sharing\open{15} & Automatic when good \\
Pen \Otwo{} measurement lost & no selected Orbit sensor with good quality & OIC01 output frozen indefinitely\open{14}; lines keep their share & Automatic, bumpless \\
One Orbit sensor fails & Modbus timeout / bad quality & Sensor drops out of the mean; CommFail alarm after delay & Automatic \\
Inlet pump lost & inlet permit FALSE (not REMOTE or heartbeat lost) & Line stops: valve closed, oxy pump stopped; other lines take over the share at once & Automatic: while oxygenation is on and the line enabled, the line starts again when its inlet pump is back in REMOTE \\
Inlet pump stopped by the pump count or rotation & inlet permit FALSE & As inlet pump lost. A line whose inlet pump is started by the pump count or rotation starts when that pump reaches REMOTE, if it is enabled & Automatic while oxygenation is on\open{23} \\
RIO link lost (InletRIO$n$ or CommonRIO) & Modbus TCP \var{CommOK} FALSE; alarm RIO comm.\ error & \emph{None yet}: PT4$n$, PT4($n$+3) (InletRIO$n$) and FT4$n$, PT40 (CommonRIO) keep their last value and good quality, so PIC4$n$, the valve permit and FIC4$n$ run on frozen values\open{28} & -- \\
\Otwo{} supply lost & PT40 (displayed only) & \emph{None}: no alarm, no interlock\open{4}\textsuperscript{,}\open{10}; FIC opens valve fully, flow stays low & -- \\
Water back into injector / supply & -- & Mechanical only (check valves)\open{6} & -- \\
Emergency stop & not wired\open{22} & (Drive quick stop is alarmed as EmergencyStop) & -- \\
\end{longtable}
\end{small}

\subsection{Communication}
\begin{itemize}[nosep]
  \item \textbf{RIO}: Modbus TCP, FC23 every 100\,ms; InletRIO$n$ read 26/write 21 words, CommonRIO read
        41/write 25 words (to be checked against the couplers, [PID] A1/A8).
  \item \textbf{Drive}: Modbus TCP, unit id 0 (\TBD\ for the real drive\open{19}), channel 0 FC23 every 100\,ms
        (write 8501--8502, read 3201--3211), channel 1 FC3 every 1\,s (read 7121).
  \item \textbf{Orbit 864}: Modbus RTU 9600\,8N1 over the 750-652, one sensor at a time, \Otwo{} block
        every second round; stay-awake register written so the sensors do not sleep.
  \item \textbf{Framo}: OPC UA; heartbeat timeout 10\,s.
  \item \textbf{HMI}: OPC UA to \var{GVLMain}.
\end{itemize}

% =============================================================================================
\section{Non-functional requirements}
\label{sec:nfreq}

\subsection{Control laws and start tuning}
All three controllers are PI (\var{PIDController}, $T_d = 0$), with output limits and anti-windup.
The values are start values from \var{ConfigValues} [DB], to be tuned at commissioning.
\begin{center}
\begin{small}
\begin{tabularx}{\textwidth}{@{}P{20mm}P{26mm}P{26mm}P{26mm}L@{}}
\toprule
Loop & PV / SP & Output & $K_p$, $T_i$ & Basis \\
\midrule
OIC01 & pen \Otwo{} mg/L; SP 8.0\,mg/L\open{13} & total \Otwo{} 0..$\sum$ kg/h & 50 (kg/h)/(mg/L), 1800\,s & Pen time constant about 44\,min, gain about 0.024 mg/L per kg/h [DB] \\
FIC4$n$ & \Otwo{} flow kg/h; SP from FY01 & valve 0--100\,\% & 0.3\,\%/(kg/h), 5\,s & Valve and meter lag about 0.2\,s; closed loop about 2\,s \\
PIC4$n$ & PT4$n$ bar; SP 4.0 & speed 20--50\,Hz & 2\,Hz/bar, 5\,s & Pump affinity law (simulator) \\
\bottomrule
\end{tabularx}
\end{small}
\end{center}
Deviation from [FD] table~13/14: [FD] has OIC01 output 0--100\,\%, three parameter sets (P 0.5/1/2,
I 600\,s) by number of pumps, and FIC in l/min with P 0.5, I 600\,s. The PLC works in kg/h with one
set; the reason is given in Section~\ref{sec:func} and [DB]\open{16}.

\subsection{Timing and resolution}
\begin{itemize}[nosep]
  \item PLC task \var{Task}: cyclic, 20\,ms. Filters and the valve ramp use the task interval read
        from the runtime.
  \item RIO and drive I/O: 100\,ms; drive last fault: 1\,s. Orbit sensors: each sensor about every
        few seconds (one RS-485 line); sensor response $T_{90} < 45$\,s [ORB].
  \item 750-455 AI: 12-bit (2048 steps over 4--20\,mA): 0--10\,bar $\rightarrow$ about 5\,mbar,
        0--150\,kg/h $\rightarrow$ about 0.07\,kg/h. 750-563 AO: 16-bit word.
  \item Drive: speed 0.1\,Hz, current 0.01\,A as read by the PLC (unit to be confirmed for the
        ATV650D15, [EAV] says ``refer to programming manual''), thermal state and power 1\,\%.
  \item Orbit: \Otwo{} 0--60\,mg/L, resolution 0.01\,mg/L, accuracy $\pm0.1$\,mg/L (0--20 mg/L) [ORB].
\end{itemize}

\subsection{Alarms}
\label{sec:alarms}
Alarms use the common \var{AlarmManagement} block. ``auto'' alarms clear themselves when the
condition is gone; ``ack'' alarms must also be acknowledged. Drive alarms are FALSE without
communication, except the communication alarm. LFT keeps the last fault after a reset, so fault
codes count only while the drive is in fault.
\begin{small}
\begin{longtable}{@{}P{38mm}P{80mm}P{10mm}P{22mm}@{}}
\caption{Alarms of the oxygen system. Codes per [EAV] (ATV600 list); ``FD/AL'' = required by [FD]/[AL].}\label{tab:alarms}\\
\toprule
Alarm (\texttt{OxySystem[i].}) & Condition & Ack & Status \\
\midrule
\endfirsthead
\toprule
Alarm & Condition & Ack & Status \\
\midrule
\endhead
\bottomrule
\endfoot
\multicolumn{4}{@{}l}{\textit{Side-stream pump drive (\texttt{OxyPump.Alarm.*}), per line}} \\
DriveCommunicationError & drive configured and Modbus TCP link not OK & auto & PLC \\
DriveError & ETA state fault; code in \var{OxyPump.Input.LastFault} & auto & PLC \\
DriveNotInRemote & ETA bit 9 = 0 (command from display terminal / forced local) & auto & PLC \\
EmergencyStop & DRIVECOM ``quick stop active'' & auto & PLC \\
DCUndervoltageError & ETA bit 4 = 0 (no mains to power part), or fault PHF (21) / USF (22) & auto & PLC \\
MotorTemperatureHigh & fault OLF (17), TH2F (110), TH3F (112), TH4F (120), TH5F (122) & auto & PLC \\
DriveTemperatureHigh & THD $\geq$ 100\,\% (\var{DriveThermalWarning_pct}), or fault OHF (16) / TJF (54) & auto & PLC \\
DriveWarning & ETA bit 7 (drive warning) & auto & PLC \\
PumpCommonError & any of the above except DriveWarning; [AL] ``Feil på OxyPumpe $n$'' & ack & PLC, FD/AL \\
StartUpTimeout & Run commanded, not turning after \var{OxyStartupTime} 10\,s & ack & PLC \\
StopTimeout & stop commanded, still turning after \var{OxyShutdownTime} 10\,s & ack & PLC \\
\multicolumn{4}{@{}l}{\textit{Line}} \\
ValvePressureLow (PAL4$n$) & oxygenation active but no valve permit: valve held closed & ack & PLC \\
Sensor quality PT4$n$, PT4($n$+3), FT4$n$ & 750-455 diagnostic bad for \var{QualityDelay} & ack & PLC \\
Low $\Delta p$ over injector & PT4($n$+3) -- PT4$n$ too small (limit from HMI) & -- & FD/AL, not in PLC\open{10} \\
Pump leakage / temp. & XS4$n$ / TSH4$n$ & -- & not in PLC\open{24} \\
\multicolumn{4}{@{}l}{\textit{Pen and supply}} \\
OIC01 deviation & $|$PV -- SP$|$ $>$ limit while in auto, with delay (HMI) & -- & FD/AL, not in PLC\open{11} \\
OT$nn$ low / high \Otwo{} & per sensor, limits and delay from HMI, only for active sensors & -- & FD/AL; L/H in DUT, not configured\open{12} \\
\Otwo{} sensor disagreement & spread between OT01..03 $>$ limit & -- & FD/AL, not in PLC\open{12} \\
\Otwo{} sensor fault & Orbit communication failure (CommFail) after delay & ack & PLC \\
\Otwo{} high-high shut-off & supersaturation (\Otwo{} HH / TDGP) closes oxygenation & -- & FD (HOLD), not in PLC\open{9} \\
PAL40 \Otwo{} not available & PT40 low & -- & FD/AL, not in PLC\open{4} \\
\end{longtable}
\end{small}

\subsection{Parameters}
\label{sec:params}
Retained in \var{GVLRetain}; defaults from \var{PFC200_Sim} \var{ConfigValues} (index $i$ = line $-1$).
\begin{small}
\begin{longtable}{@{}P{62mm}P{30mm}P{62mm}@{}}
\caption{Configuration parameters of the oxygen system.}\label{tab:params}\\
\toprule
Parameter & Default & Meaning \\
\midrule
\endfirsthead
\toprule
Parameter & Default & Meaning \\
\midrule
\endhead
\bottomrule
\endfoot
\var{PenConfig.OxySetpoint} & 8.0 mg/L & OIC01 setpoint (placeholder)\open{13} \\
\var{PenConfig.OxyController} & $K_p$ 50, $T_i$ 1800 s, Min 0, Max 0 & OIC01; Max 0 = what the lines can deliver \\
\var{Orbit864SensorConfig[k]} & slaves 1--6; EnableForAutoControl for 1--3 & sensor address, polling, selection for OIC01 \\
\var{OxySystemConfig[i].OxyPump.MinSpeed_Hz / MaxSpeed_Hz} & 20 / 50 Hz & PIC output range\open{19} \\
\var{.OxyPump.PresetSpeed_Hz} & 40 Hz & speed while starting, with bad PT, or PIC disabled \\
\var{.OxyPump.RunDelay} & 30 s & pump running before oxygen may be added \\
\var{.OxyPump.DriveThermalWarning_pct} & 100 \% & THD limit for DriveTemperatureHigh \\
\var{.OxyStartupTime / .OxyShutdownTime} & 10 s / 10 s & start / stop timeout alarms \\
\var{.PressureController} & Enabled, $K_p$ 2 Hz/bar, $T_i$ 5 s & PIC4$n$ \\
\var{.PressureSetpoint} & 4.0 bar & PIC4$n$ setpoint (abs. assumed)\open{17}\textsuperscript{,}\open{18} \\
\var{.ValveOpenPressure} & 3.0 bar & valve permit limit, same scale as setpoint \\
\var{.FlowController} & Enabled, $K_p$ 0.3 \%/(kg/h), $T_i$ 5 s & FIC4$n$ \\
\var{.FlowMax} & 104.2 kg/h & most \Otwo{} flow of one line\open{7} \\
\var{.ValveRampTime} & 2 s & valve command 0--100 \% at most in this time \\
\var{OxyPumpPressureConfig[i]} & 0--10 bar, $f_c$ 1 Hz & PT4$n$ scaling, filter, quality \\
\var{OxynatorOutletPressureConfig[i]} & 0--10 bar, $f_c$ 1 Hz & PT4($n$+3) \\
\var{OxygenFlowConfig[i]} & 0--150 kg/h, $f_c$ 1 Hz & FT4$n$\open{3} \\
\var{InletRIOConfig[i].DriveIP / DrivePort / DriveUnitId} & sim 10.1.2.101 / 5131--5133 / 0 & drive address\open{19} \\
\end{longtable}
\end{small}

\subsection{Instruments}
% casfd's instrumenttable columns add up to 160 mm; with the default column separation the table
% is wider than the text, so the separation is reduced for this table.
\setlength{\tabcolsep}{1.4pt}
\begin{instrumenttable}{Instruments and actuators of the oxygen system. TBD = not yet known.}
PT40 & \Otwo{} supply pressure & Normex panel header & 4--20 mA, CommonRIO -R1-008 ch4 & \TBD{} (PLC default 0--10 bar) & 12 bit & display\open{4} \\
FT41..43 & \Otwo{} mass flow per line (FM-01.$n$) & Normex panel & 4--20 mA, CommonRIO -R1-008 ch1--3 & 0--150 kg/h in PLC, \TBD & 12 bit; $\pm$1.5\,\% m.v. & FIC PV\open{3} \\
OV41..43 & \Otwo{} control valve (MV-01.$n$) & Normex panel & 4--20 mA, CommonRIO -R1-014/015 & 0--100 \% & 16 bit & FIC out\open{2} \\
PT41..43 & water pressure after side-stream pump & side-stream line & 4--20 mA, InletRIO$n$ ch3 & 0--10 bar (PLC), \TBD & 12 bit & PIC PV, permit \\
PT44..46 & water pressure after injector & side-stream line & 4--20 mA, InletRIO$n$ ch4 & 0--10 bar (PLC), \TBD & 12 bit & display \\
PT-01.1..3 & \Otwo{} pressure after valve & Normex panel & \TBD{} (not connected) & $-1$..+6 bar & -- & \TBD \\
OT01..03 & \Otwo{}, temp., salinity, inside pen & pen, near outlet & Modbus RTU, 750-652 & 0--60 mg/L & 0.01 mg/L & OIC PV \\
OT21..23 & \Otwo{}, temp., salinity, outside & 3/30/50 m from pen & Modbus RTU, 750-652 & 0--60 mg/L & 0.01 mg/L & reference \\
SC41..43 & side-stream pump drive & DC$n$ & Modbus TCP & 20--50 Hz & 0.1 Hz & PIC out \\
XS41, TSH41 & pump leakage, motor temp. & DC1 MiniCAS II & DI InletRIO1 5.7/5.8 & on/off & -- & not used\open{24} \\
PT11, QT11 & TDGP, CO\textsubscript{2} (Solu-Blu) & pen near outlet & 4--20 mA, CommonRIO -R1-007 & \TBD & 12 bit & display; HH?\open{9} \\
\end{instrumenttable}
\setlength{\tabcolsep}{6pt}

\subsection{Drive Modbus registers}
Checked against the ATV600 communication parameter list [EAV] (logic address = Modbus register).
\begin{center}
\begin{small}
\begin{tabularx}{\textwidth}{@{}P{12mm}P{12mm}P{34mm}P{12mm}P{22mm}L@{}}
\toprule
Code & Reg. & Name & R/W & Unit & PLC use \\
\midrule
CMD & 8501 & Command register & W & bits & DRIVECOM control word, FC23 every 100\,ms \\
LFR & 8502 & Reference frequency & W & 0.1 Hz & speed setpoint \\
ETA & 3201 & Status register & R & bits & state, fault, remote (bit 9), warning (bit 7), mains (bit 4) \\
RFR & 3202 & Motor frequency & R & 0.1 Hz & actual speed, Running \\
LCR & 3204 & Motor current & R & see [EAV] & current (PLC: 0.01 A) \\
THD & 3209 & Drive thermal state & R & 1 \% & DriveTemperatureHigh \\
OPR & 3211 & Motor power in \% & R & 1 \% & display \\
LFT & 7121 & Last error occurred & R & enum & fault code, FC3 every 1\,s \\
ETHL & 7021 & Ethernet error response & cfg & enum & set in drive: factory freewheel stop; options include fallback speed LFF (7080) and speed maintain\open{19} \\
\bottomrule
\end{tabularx}
\end{small}
\end{center}
All register numbers and the fault codes used by \var{OxyPumpATV650} (16 OHF, 17 OLF, 21 PHF,
22 USF, 54 TJF, 110/112/120/122 TH2F--TH5F, communication fault ETHF 150) match [EAV]. The program
design [PD] still documents the ATV320 block with codes from the ATV320 list NVE41316 (including
PTFL 49 and OTFL 50, which the ATV600 list does not have); the codes shall be re-checked on the
real drive's firmware at commissioning and [PD] updated\open{21}.

\subsection{HMI operator functions}
The production HMI is not yet specified\open{25}; the present CAS HMI is a read-only overview of
pen \Otwo, level and valve signals. The PLC offers the following operator functions in
\var{GVLMain}; [FD] requires them on the HMI.
\begin{small}
\begin{longtable}{@{}P{48mm}P{60mm}P{46mm}@{}}
\caption{Operator functions.}\label{tab:hmi}\\
\toprule
Function & PLC variable & [FD] \\
\midrule
\endfirsthead
\toprule
Function & PLC variable & [FD] \\
\midrule
\endhead
\bottomrule
\endfoot
Oxygen Start / Stop (pen) & \var{Pen.Command.OxygenStart}, \var{.OxygenStop}; \var{Pen.Status.OxygenOn} & pulses; enabled lines follow their inlet pumps \\
Enable / disable oxygen line $n$ & \var{OxyConfig[i].Enable} & retained \\
Pump auto / manual, start, stop, speed, reset & \var{OxySystem[i].OxyPump.Command.Auto/Manual/Start/Stop/Setpoint/Reset} & manual start allowed \\
Valve auto / manual, manual position & \var{OxySystem[i].OxyValve.Command.Manual/Auto} (pulses; mode in \var{.Status.Manual}), \var{.ManualSetpoint} & -- \\
OIC01 setpoint & \var{PenConfig.OxySetpoint} & \% or mg/L?\open{13} \\
OIC01 manual output & \var{OxyController.Command.Manual/ManualOutput} & -- \\
Select sensors for control & \var{Orbit864SensorConfig[k].EnableForAutoControl} & selector; inactive sensors yellow \\
Preset speed, run delay, tuning & \var{OxySystemConfig[i]}, \var{PenConfig} & preset speed from HMI \\
Drive info panel & \var{OxyPump.Input}: ready, remote, running, power, warning, error, speed, current, LFT & table 12 \\
Alarm limits (deviation, L/H, $\Delta p$) & not yet in PLC & tables 15, 17\open{11}\textsuperscript{,}\open{12} \\
\end{longtable}
\end{small}

\subsection{Deviations from the Norwegian FD}
\begin{small}
\begin{longtable}{@{}P{8mm}P{58mm}P{58mm}P{28mm}@{}}
\caption{Implemented behaviour that differs from [FD] rev~B.}\label{tab:dev}\\
\toprule
No. & [FD] & PLC 2026-10-21 & Action \\
\midrule
\endfirsthead
\toprule
No. & [FD] & PLC 2026-10-21 & Action \\
\midrule
\endhead
\bottomrule
\endfoot
1 & OIC01 output 0--100\,\% to the valves, or cascade to FIC as an alternative & cascade: OIC01 total kg/h shared over FIC4$n$ & update FD\open{16} \\
2 & three OIC parameter sets by number of oxy pumps & one set (total flow output makes the gain independent) & update FD \\
3 & FIC SP in l/min, scaling 0--20, P 0.5, I 600\,s & kg/h, 0--150, $K_p$ 0.3, $T_i$ 5\,s & update FD; confirm unit\open{3} \\
4 & PV = one selected sensor (OT03), switch to lowest on fault & mean of the selected good sensors & decide \\
5 & valve: PWM solenoid ``via 4--20\,mA'' & proportional 4--20\,mA, 2\,s ramp in PLC & Normex\open{2} \\
6 & \Otwo{} fault: freeze output (``HOLD'', sentence unfinished) & freeze indefinitely & decide\open{14} \\
\end{longtable}
\end{small}

% =============================================================================================
\section{Open issues and decisions}
\label{sec:open}

\textbf{Decisions.}
\begin{description}[nosep, leftmargin=12mm, style=sameline]
  \item[DEC1] 2026-10-21: oxygen supply and distribution panel from Normex.
  \item[DEC2] Side-stream pump drive ATV650D15N4 on Modbus TCP (embedded Ethernet).
\end{description}

\textbf{Open issues.} Questions for Normex are marked (N).
\begin{small}
\begin{longtable}{@{}P{9mm}P{98mm}P{48mm}@{}}
\toprule
No. & Issue & Proposed action / owner \\
\midrule
\endfirsthead
\toprule
No. & Issue & Proposed action / owner \\
\midrule
\endhead
\bottomrule
\endfoot
OI1 & (N) Supply type (LOX tank, generator, or customer supply line) and normal / minimum supply pressure (sketch: ``2.5 bar ?''; Linde sheet: 10 bar). & Normex to state; sets PT40 range and PAL40 limit. \\
OI2 & (N) Control valve type: [FD] PWM solenoid, PLC proportional 4--20\,mA with 2\,s ramp, Normex sketch motorised valve MV-01.$n$. Signal, stroke time, fail position (closed on signal loss?) and position feedback. & Normex; adjust \var{ValveRampTime} and FIC tuning to the stroke time. \\
OI3 & (N) Flow meter unit (kg/h vs Nl/min / Nm\textsuperscript{3}/h), range for \Otwo, signal (4--20\,mA non-isolated or Modbus RTU). PLC assumes 4--20\,mA = 0--150 kg/h. & Normex; update \var{OxygenFlowConfig}. \\
OI4 & (N) Low supply pressure signal: which transmitter (PT40 / PI-01.4) and limit; PAL40 ``oxygen not available'' and a supply pressure interlock are not in the PLC (PT40 only displayed, misnamed \var{OxygenInjectionPressure}). & Add PAL40 and interlock; rename. \\
OI5 & (N) PLC-controlled master shut-off valve: \emph{decided 2026-10-22}: no on/off solenoid or shut-off valves upstream of the control valves (budget). The control valve is the only shut-off; the PLC holds it at 4\,mA without the valve permit (pump running, pressure after the pump above ar{ValveOpenPressure}), the inlet permit or on emergency stop. \Otwo{} HH shut-off can only close the control valves. & Closed (register OX-17). \\
OI6 & (N) Non-return protection: HAZID ``oxygen pressure lost, water enters injector and \Otwo{} tank''. Check valve per line (Linde had CV1--3; Normex sketch unclear: OV-02.$n$). & Normex to confirm. \\
OI7 & (N) Injector model and gas ratio: \var{FlowMax} 104.2 kg/h = CT10 at 15\,\% v/v; 34.7 kg/h at 5\,\%. Which applies in sea water at Heimdall? & Chucaotech / Normex; set \var{FlowMax}. \\
OI8 & (N) Signal list and interface FD for the Normex panel (terminals, 24\,V supply, cable list). & Normex. \\
OI9 & \Otwo{} high-high / supersaturation (TDGP) shut-off asked by [FD] (HOLD); not in the PLC. Needs limits and an actuator (OI5). & Decide limits; implement. \\
OI10 & Injector low differential pressure alarm ([FD] table 15): \var{OxynatorOutletPressure} is read but unused. & Implement with HMI limit. \\
OI11 & OIC01 deviation alarm ([FD] table 15): not implemented. & Implement. \\
OI12 & \Otwo{} sensor disagreement alarm and per-sensor low/high alarms: not implemented (L/H fields exist in the Orbit DUT, not configured). & Implement; set limits. \\
OI13 & \Otwo{} setpoint 8.0\,mg/L is a placeholder; [FD] asks \% or mg/L. [O2S] uses \% saturation: $\geq$80\,\%, 90\,\% at 16\degC. A \% setpoint needs temperature/salinity correction. & Customer / biology to decide. \\
OI14 & Loss of pen \Otwo{} measurement: OIC01 holds its output indefinitely ([FD] sentence unfinished). Hold, fixed fallback flow, or time limit? & Decide; implement. \\
OI15 & Loss of line flow measurement: valve holds last position, line flow counted as 0, so the other lines may over-dose. Hold, close, or fixed opening? & Decide; implement. \\
OI16 & Units and tuning differ from [FD] (Table~\ref{tab:dev}); deliberate. & Accept; update [FD] to rev~C. \\
OI17 & Pressure transmitters: make, range, absolute or gauge unknown ([SO] ``find sensors''). \var{PressureSetpoint} 4.0 and \var{ValveOpenPressure} 3.0 must change together. Transmitter before the pump wanted? & Select transmitters; set ranges. \\
OI18 & Pump head vs. pressure setpoint: NX 3153 gives 10.2\,m at 5.05\,m\textsuperscript{3}/min, 18.3\,m shut-off at 50\,Hz; the simulator assumes 4.5\,bar rise. Required injector inlet pressure from Chucaotech? & Check suction pressure and injector requirement; reset SP. \\
OI19 & ATV650: IP address and unit id, comm-loss behaviour (ETHL; [FD] expects running on at preset speed), 20--50\,Hz range and 40\,Hz preset vs. pump curve. & Set at commissioning; decide ETHL (LFF or RLS). \\
OI20 & Drive cabinet 874.004 has no Ethernet drawn; RJ45 DIN socket proposed. & Electrical design. \\
OI21 & Drive fault codes: PLC uses ATV600 codes (match [EAV]); [PD] still lists ATV320 (NVE41316). & Verify on the real drive; update [PD], [DB]. \\
OI22 & Emergency stop not wired: \var{OxySystem[i].Input.EmergencyStop} never written; [FD] wants E-stop to stop PU4$n$ and close OV4$n$. & Wire (Framo E-stop or hardwired) ([PID] A3). \\
OI23 & Inlet permit: PLC uses inlet pump REMOTE + live heartbeat; [FD] says ``inlet pump running''. Auto-restart of oxygen lines after power return ([FD] table 2.2). & Confirm REMOTE implies running; decide restart. \\
OI24 & Pump leakage / temperature (XS, TSH) only on RIO1, not used in logic. & Add inputs on RIO2/3 and alarms ([PID] B8). \\
OI25 & Production HMI not specified; operator functions of Table~\ref{tab:hmi}. & HMI design. \\
OI26 & Oxygen capacity with one inlet pump out: two lines carry only about 560--730\,t [DB]; fallback is emergency oxygenation and stopped feeding. & Project / customer. \\
OI27 & PT numbering [FD] (PT41/43/45 in, PT42/44/46 out) vs. drawings (PT41--43 after pump, PT44--46 after injector). & Update [FD]. \\
OI28 & RIO link loss freezes the line's analog values with good quality (register OX-28; the inlet level pairs got this fix in CAS-FD-10 rev~B). & Mark the RIO's sensors bad while \var{CommOK} is FALSE; line leaves auto, valve per OI15. \\
\end{longtable}
\end{small}

\end{document}
